\bigskip
\noindent Autres pistes de l'onglet Findings pour ce thème, dans des cotes qui n'ont pas de section ici :\par

\begin{encadre}
\begin{theoreme*}[Grothendieck, pages 10 et 11, dualité locale en dimension un pour tous les modules\piste{101-sheaf-duality-dim-one}]\nonverifie
Sur $X = \operatorname{Spec} A$, où $A$ est un anneau de valuation discrète complet de corps des fractions $K$, les applications de dualité locale $\mathrm{Ext}^{1-i}(X; F, \mathcal O_X) \to \mathrm{Hom}_A(H^i_x(F), K/A)$ sont des isomorphismes pour tout $\mathcal O_X$-module $F$, quasi-cohérent ou non. La démonstration utilise la complétude de $A$ par $\mathrm{Ext}^1_A(K, A) = 0$.
\end{theoreme*}
\end{encadre}
{\small La fin de l'argument d'injectivité de $\varphi^1$ n'est lisible qu'en partie, et la lecture le restitue. Mathlib n'a pas la cohomologie locale des faisceaux de modules non quasi-cohérents. On se reportera aux notes de Grothendieck (\lot{101}{1}{10}).\par}

\begin{encadre}
\begin{contrexemple*}[pages 10, 14 et 15, dualité locale en dimension deux\piste{101-sheaf-duality-dim-two-fails}]\nonverifie
La première phrase de la page~10 affirme les isomorphismes de dualité locale en toute dimension ; ils tombent en défaut en dimension~$2$. Soient $A$ un anneau local régulier complet de dimension~$2$, $f \in \mathfrak m$ non nul, $U = X_f$ et $F = j_!\mathcal O_U$ sur $X = \operatorname{Spec} A$. L'application $A_f = \mathrm{Ext}^0(X; F, \mathcal O_X) \to \mathrm{Hom}_A(H^2_x(F), E(k))$ est injective et non surjective : la dualité locale à duaux de Matlis algébriques ne vaut pas pour les $\mathcal O_X$-modules non quasi-cohérents.
\end{contrexemple*}
\end{encadre}
{\small La page~14 conclut \q{$\varphi$ n'est pas surjective}, souligné deux fois, sans démonstration, et la page~15 dit qu'il faut des duaux topologiques. Mathlib n'a pas la cohomologie locale ni la dualité de Matlis. On se reportera aux notes de Grothendieck (\lot{101}{1}{10}).\par}

\begin{encadre}
\begin{theoreme*}[Grothendieck, pages 10 et 11, noyau de l'image réciproque et indice\piste{32-pullback-kernel-index}]\nonverifie
Soient $f : X \to Y$ propre sur un corps $k$, $Y$ irréductible et lisse de point générique $\eta$, $X_\eta$ non vide, et $N$ le pgcd des degrés des $0$-cycles de $X_\eta$. Il existe une application $H^*(X, G_X) \to H^*(Y, G)$, fonctorielle en le complexe de coefficients étale $G$, dont le composé avec $f^*$ est $N \cdot \mathrm{id}$. Le noyau de $f^*$ est donc tué par l'indice de la fibre générique, sans hypothèse sur la dimension relative.
\end{theoreme*}
\end{encadre}
{\small Le calcul de la page~10 suppose le formalisme de Gysin du \S\,4 pour des $Z_i$ éventuellement singuliers, que le dossier énonce sans le démontrer. Mathlib n'a pas la cohomologie étale. On se reportera aux notes de Grothendieck (\lot{32}{1}{10}).\par}

\begin{encadre}
\begin{theoreme*}[Grothendieck, pages 165 à 169, $\mathrm{Ext}^1(\mathbf G_a, \mathbf Z/p\mathbf Z)$ et perfection\piste{43-ext-ga-zp-perfection}]\nonverifie
Sur une base quelconque $S$ de caractéristique~$p$, $\mathrm{Ext}^1_{S\text{-gr}}(\mathbf G_a, \mathbf Z/p\mathbf Z) \simeq \Gamma(S, \mathcal O^{p^{-\infty}})$, sections globales de la perfection $\varinjlim(\mathcal O \to \mathcal O \to \cdots)$ pour $x \mapsto x^p$. Le foncteur $\mathrm{Ext}^1(\mathbf G_a, \mathbf Z/p\mathbf Z)$ est donc représentable parmi les pro-schémas parfaits sur~$S$.
\end{theoreme*}
\end{encadre}
{\small Les feuillets en esquissent une démonstration incomplète (pages 166 à 168) : ils démontrent le lemme qui identifie $\mathrm{Ext}^1$ aux classes primitives de $H^1(G, \Gamma)$ et son invariance par homéomorphisme universel, mais pas la formule elle-même, et le mot \q{parfaits} est incertain. Mathlib n'a pas les $\mathrm{Ext}$ de schémas en groupes. On se reportera aux notes de Grothendieck (\lot{43}{9}{165}).\par}

\begin{encadre}
\begin{theoreme*}[D'après les pages 70 à 74, forme corrigée par l'édition, concentration de la cohomologie locale\piste{43-local-cohomology-concentration}]\nonverifie
Soient $A$ un anneau local régulier de dimension $d$ contenant un corps algébriquement clos $k$, et $G$ un groupe algébrique commutatif sur $k$. La cohomologie locale de Zariski $H^i_{\mathfrak a}(\operatorname{Spec} A, G)$ au point fermé est nulle pour $i \neq d$ ; elle vaut $G(K)/G(A)$ pour $d = 1$ et $H^{d-1}(\operatorname{Spec} A - \mathfrak a, U)$ pour $d \geq 2$, où $U$ est la partie unipotente de $G$. Le foncteur $G \mapsto H^d_{\mathfrak a}(G)$ est donc exact sur les suites localement triviales, ce qui sert à définir la jacobienne locale $J_{A/k}$ de la page~74.
\end{theoreme*}
\end{encadre}
{\small La clôture algébrique de $k$ est ajoutée par l'édition, et presque tous les mots de la prose sont de lecture incertaine. Mathlib n'a pas la cohomologie locale à valeurs dans un groupe algébrique. On se reportera aux notes de Grothendieck (\lot{43}{4}{70}).\par}

\begin{encadre}
\begin{theoreme*}[Grothendieck, pages 40 à 43, permanence de la condition de Lefschetz\piste{58-lef-permanence-projective-flat}]\nonverifie
Soient $X$ quasi-projectif sur un anneau noethérien et $Y$ un fermé de $X$ tels que la paire $(X, Y)$ vérifie $(\mathrm{LG})_n$ : $H^i(X, F) \to H^i(X_{/Y}, F_{/Y})$ est bijectif pour $i \leq n$ et injectif pour $i = n+1$ pour tout $F$ localement libre. Alors pour tout $f : W \to X$ projectif et plat, la paire $(W, f^{-1}(Y))$ vérifie $(\mathrm{LG})_n$ ; de même pour la condition $\mathrm{Lef}$ de SGA~2, X. Rien n'est demandé aux fibres de~$f$.
\end{theoreme*}
\end{encadre}
{\small Démonstration donnée par l'édition, non vérifiée en Lean.\par}
{\small SGA~2 (exposés X et XII) obtient ces conditions par des hypothèses de profondeur sur les fibres. Mathlib n'a pas les complétions formelles ni leur cohomologie. On se reportera aux notes de Grothendieck (\lot{58}{2}{40}).\par}

\begin{encadre}
\begin{theoreme*}[Grothendieck, pages 45 à 51, comparaison des Hom hors de $Y$\piste{58-local-hom-comparison-depth-off-y}]\nonverifie
Soient $A$ un anneau local noethérien, quotient d'un anneau régulier et $t$-adiquement complet, $X'$ son spectre épointé et $Y' = V(t) \cap X'$. Pour $F$ et $G$ cohérents sur $X'$, $\mathrm{Hom}(F, G) \simeq \mathrm{Hom}(\hat F, \hat G)$ le long de $Y'$ dès que $G$ est de profondeur $\geq 2$ aux points fermés de $X' - Y'$. Aucune profondeur n'est demandée aux points de $Y'$, alors que SGA~2, X, exemple~2.1~(i), la demande en tout point fermé de~$X'$.
\end{theoreme*}
\end{encadre}
{\small Les feuillets en esquissent une démonstration incomplète (pages 45 à 51) : la page~47 marque d'un \q{?} le pas $t$-régulier et s'arrête. Mathlib n'a pas les faisceaux cohérents sur un spectre épointé ni leur complétion formelle. On se reportera aux notes de Grothendieck (\lot{58}{3}{45}).\par}
